%% macros.tex —— 正文使用的排版宏
%%
%% 源码路径只出现在 \srcfile / \fullsource 的参数里，正文中的代码一律由 minted
%% 从 ../include/acm/ 读取。表格用 longtable（begin/end 都在宏里，不做向前扫描）。

% ---------------------------------------------------------------- 行内代码
% \code{add_edge}：先整体 \detokenize，再逐字符输出，在 _ . : / , ( - 处允许断行。
\ExplSyntaxOn
\tl_new:N \l__acm_code_tl
\NewDocumentCommand \code { m }
    {
        \group_begin:
        \tl_set:Nx \l__acm_code_tl { \detokenize {#1} }
        \tl_replace_all:Nnn \l__acm_code_tl { ~ } { \c_space_tl }
        \texttt
            {
                \tl_map_inline:Nn \l__acm_code_tl
                    {
                        ##1
                        \str_case_e:nn { \tl_to_str:n {##1} }
                            {
                                {_} { \allowbreak }
                                {.} { \allowbreak }
                                {:} { \allowbreak }
                                {/} { \allowbreak }
                                {,} { \allowbreak }
                                {(} { \allowbreak }
                                {-} { \allowbreak }
                            }
                            { }
                    }
            }
        \group_end:
    }
\ExplSyntaxOff

% 头文件路径（相对仓库根）
\newcommand{\srcfile}[1]{\code{include/acm/#1}}

% ---------------------------------------------------------------- 接口表
% \apitable{ \api{签名}{说明} ... }
\newcommand{\apitable}[1]{%
    \par\medskip\noindent
    \begingroup
    \small
    \renewcommand{\arraystretch}{1.38}%
    \begin{longtable}{@{}>{\raggedright\arraybackslash}p{0.40\linewidth}>{\raggedright\arraybackslash}p{\dimexpr0.60\linewidth-2\tabcolsep\relax}@{}}
        \toprule
        \textbf{接口} & \textbf{说明} \\
        \midrule
        \endfirsthead
        \toprule
        \textbf{接口} & \textbf{说明} \\
        \midrule
        \endhead
        \bottomrule
        \endfoot
        #1%
    \end{longtable}%
    \endgroup
    \par\medskip
}

\newcommand{\api}[2]{\code{#1} & #2 \\}

% ---------------------------------------------------------------- 完整源码
% \fullsource{max_flow.hpp}：与其它小节统一的标题，不重复源码路径
\newcommand{\fullsource}[1]{%
    \par\needspace{6\baselineskip}%
    \subsection{完整源码}%
    \tcbinputlisting{acmcode, listing file = {../include/acm/#1}}%
}

% ---------------------------------------------------------------- 速查索引
% \quickrefrow{模块}{头文件}{复杂度}{sec:label}
\newcommand{\quickrefrow}[4]{#1 & #2 & #3 & \pageref{#4} \\}
